\chapter{Call the Catalog API}

The store can reach the catalog jail at the network layer, but its application
code does not yet make that request. We will introduce the integration through
a small client object. Calling it directly before adding a controller keeps
the failure boundary narrow: any problem belongs to HTTP, JSON, or the catalog
service rather than HTML rendering.

\section{Exercise 13A: Write the catalog client}

In the store jail, work as the application user:

\begin{lstlisting}[language=sh]
su -l store
cd /usr/local/www/store
mkdir -p app/services
vi app/services/catalog_client.rb
\end{lstlisting}

Add the following class:

\begin{lstlisting}[language=Ruby]
require "json"
require "net/http"
require "uri"

class CatalogClient
  class Error < StandardError; end

  DEFAULT_BASE_URL = "http://10.17.89.24:3000"

  def initialize(base_url: ENV.fetch(
    "CATALOG_URL", DEFAULT_BASE_URL
  ))
    @base_url = base_url
  end

  def products
    uri = URI.join("#{@base_url}/", "products")
    http = Net::HTTP.new(uri.host, uri.port)
    http.open_timeout = 2
    http.read_timeout = 3

    response = http.get(uri.request_uri)
    unless response.is_a?(Net::HTTPSuccess)
      raise Error, "catalog returned HTTP #{response.code}"
    end

    JSON.parse(response.body)
  rescue JSON::ParserError => error
    raise Error, "catalog returned invalid JSON: #{error.message}"
  rescue SocketError, SystemCallError,
         Net::OpenTimeout, Net::ReadTimeout => error
    raise Error, "catalog request failed: #{error.message}"
  end
end
\end{lstlisting}

The private address is a useful default for this fixed experiment. The
\texttt{CATALOG\_URL} environment variable allows tests and later deployments
to replace it without editing application code. Short timeouts also prevent a
failed dependency from leaving the store request hanging indefinitely.

\section{Exercise 13B: Probe the client directly}

Make sure the catalog Rails server is running in the catalog jail:

\begin{lstlisting}[language=sh]
su -l catalog
cd /usr/local/www/catalog
bin/rails server --binding 10.17.89.24 --port 3000
\end{lstlisting}

Leave that process running. In the store jail, invoke the new client through
the store's Rails environment:

\begin{lstlisting}[language=sh]
cd /usr/local/www/store
bin/rails runner 'pp CatalogClient.new.products'
\end{lstlisting}

The command should print an array containing the catalog product created
earlier, including SKU \texttt{BSD-001}. This proves that the store's Ruby
process can perform the HTTP request and parse the catalog's JSON response.

\section{Exercise 13C: Add the products page}

Generate a controller and view without changing the routes automatically:

\begin{lstlisting}[language=sh]
bin/rails generate controller Products index \
  --skip-routes
\end{lstlisting}

Replace \texttt{app/controllers/products\_controller.rb} with:

\begin{lstlisting}[language=Ruby]
class ProductsController < ApplicationController
  def index
    @products = CatalogClient.new.products
  rescue CatalogClient::Error => error
    @products = []
    @catalog_error = error.message
    render status: :bad_gateway
  end
end
\end{lstlisting}

Replace \texttt{app/views/products/index.html.erb} with:

\begin{lstlisting}
<h1>Products</h1>

<% if @catalog_error %>
  <p role="alert">
    Catalog unavailable: <%= @catalog_error %>
  </p>
<% elsif @products.empty? %>
  <p>No products are currently available.</p>
<% else %>
  <ul>
    <% @products.each do |product| %>
      <li>
        <strong><%= product.fetch("name") %></strong>
        (<%= product.fetch("sku") %>)
        &mdash;
        <%= number_to_currency(
          product.fetch("price_cents") / 100.0
        ) %>
      </li>
    <% end %>
  </ul>
<% end %>
\end{lstlisting}

Finally, open \texttt{config/routes.rb}. Add this line inside its existing
routes block:

\begin{lstlisting}[language=Ruby]
root "products#index"
\end{lstlisting}

The controller knows only the catalog client's interface. It renders the
returned products on success and an explicit \texttt{502 Bad Gateway} page
when the dependency reports a controlled integration error.

\section{Exercise 13D: Serve the integrated page}

Keep the catalog server running. In the store jail, start its Rails server on
the store jail's private address:

\begin{lstlisting}[language=sh]
cd /usr/local/www/store
bin/rails server --binding 10.17.89.25 --port 3000
\end{lstlisting}

Leave the process running and request the page from the FreeBSD host:

\begin{lstlisting}[language=sh]
curl -i http://10.17.89.25:3000/
\end{lstlisting}

Expect \texttt{HTTP/1.1 200 OK} and HTML containing the FreeBSD Handbook. The
request now crosses the store-to-catalog boundary while producing a
browser-facing response.

\section{Checkpoint}

Stop after the host receives the rendered product page. The store and catalog
are now integrated manually; the next chapter will turn that behavior into a
repeatable automated test.
